% Generated draft: requires mathematical and source review.
\documentclass[11pt]{article}
\usepackage[margin=25mm]{geometry}
\usepackage{fontspec}
\setmainfont{Latin Modern Roman}
\usepackage{amsmath,amssymb,mathtools,unicode-math}
\setmathfont{Latin Modern Math}
\usepackage{enumitem,needspace,xurl,hyperref}
\setlength{\parindent}{0pt}
\setlength{\parskip}{5pt}
\setlength{\emergencystretch}{3em}
\newcommand{\N}{\mathbb{N}}
\newcommand{\Z}{\mathbb{Z}}
\newcommand{\Q}{\mathbb{Q}}
\newcommand{\R}{\mathbb{R}}
\newcommand{\C}{\mathbb{C}}
\newcommand{\F}{\mathbb{F}}
\newcommand{\A}{\mathbb{A}}
\newcommand{\PP}{\mathbb{P}}
\newcommand{\E}{\mathbb{E}}
\newcommand{\eps}{\varepsilon}
\newcommand{\dd}{\,\mathrm{d}}
\begin{document}
\section*{UnsolvedMath ID: 11000228}
\textbf{Source problem code:} catalogue entry
\subsection{Definitions and mathematical statement}
Write $\mathcal Q(k_1,\ldots,k_s)$ for the relevant moduli stratum of meromorphic quadratic differentials $q$ on compact Riemann surfaces, with prescribed zero/pole orders $k_i$, where order $-1$ is a simple pole. The degree relation is $\sum_i k_i=4g-4$. Thus the listed strata
\[
 \mathcal Q(-1,9),\quad\mathcal Q(-1,3,6),\quad
 \mathcal Q(-1,3,3,3),\quad\mathcal Q(12)
\]
lie in genera $3,3,3,4$, respectively. The precise convention concerning global squares of abelian differentials and labeled singularities must agree with the source; it is not restated in the excerpt.

\textbf{Problem 21 (exceptional strata).} Find a geometric invariant distinguishing the different connected components of each of these four strata. More precisely, for each listed stratum $\mathcal Q$, seek an invariant $I_{\mathcal Q}$ of its differential/flat geometry that is constant on each connected component and takes different values on different components. Equivalently, it should induce an injective labeling of $\pi_0(\mathcal Q)$ by geometrically defined data. The problem is to distinguish components within each fixed stratum, not merely distinguish the four different singularity-order lists.

The source says that the known distinguishing invariant, the extended Rauzy class, is combinatorial rather than geometric, citing [La]. Merely renaming that class is not a geometric construction. The excerpt gives no proposed invariant, number of components, or formal criterion for what geometric description is acceptable.
\subsection{Short English statement}
Find a geometric way to tell apart connected components in the four specified exceptional quadratic-differential strata, where the source currently supplies only extended Rauzy classes as combinatorial labels.
\subsection{Sources}
\begin{itemize}
\item Local UnsolvedMath registry, permanent ID 11000228, source code catalogue entry; no proposer or proposal year supplied.
\item Farb, edited volume \emph{Problems on Mapping Class Groups and Related Topics}, Problem 21, PDF page 264; \url{https://www.math.uchicago.edu/~farb/papers/mcgbook.pdf}. The registry records PDF-font-structure extraction and access on 2026-07-29; the inventory distinguishes 329 labeled headings by source page.
\item Source label [La] for the extended-Rauzy-class component distinction; the complete bibliography entry is absent. The research summary supplies no verified later status.
\end{itemize}
\end{document}
